\begin{abstractCZ}
Práce navrhuje systém pro záznam trojrozměrné trajektorie vozidla Formula
Student vůči místní základně. Uchovává surové observace GNSS, měření IMU
a navigační řešení v reálném čase pro následné PPK zpracování a vyhlazení.
Pro rover je zvolen Septentrio mosaic-X5 na simpleRTK3B Pro, pro přenosnou
základnu u-blox ZED-F9P na simpleRTK2B Budget. Adafruit ICM-20948 nahrazuje
původně uvažovanou ISM330DHCX kvůli magnetometru a dostupnosti. Text
zohledňuje vyšší šum, magnetické rušení a ukončovanou životnost čipu.

Raspberry Pi 5 se 2~GB RAM, chlazením, M.2 HAT+ a vlastním SSD bude
s Raspberry Pi OS číst IMU přímo přes SPI a přerušení. Pro akcelerometr
a gyroskop je navrženo \SI{562.5}{Hz}, pro magnetometr nejvýše
\SI{100}{Hz}. Časovou vazbu GNSS PPS a GPIO je nutné ověřit.
Korekce RTCM má přenášet spoj LoRa v pásmu 868~MHz: dvě ze čtyř sad
XIAO ESP32S3 + Wio-SX1262, další dvě slouží pro vývoj a rezervu.
Základna má mít vlastní raw logger a webovou správu. Firmware a přenos
vyžadují ověření. Dvě antény ANN-MB2 jsou určeny pro vlastní montáž.
Oceněná elektronika stojí přibližně 24~290~Kč bez DPH, dopravy a instalace.
Navazující kapitola rozšiřuje návrh o celovozidlový DAQ: pohyb roll/heave,
otáčky kol, teploty třmenů a pneumatik. Volí centrální ADC u Raspberry,
dva CANy pro auto a posuzuje časování i napájení z 8 až 16 V. Tyto doplňky
nejsou součástí původního rozpočtu. Text je návrh, nikoli výsledek měření.
\end{abstractCZ}

\begin{keywordsCZ}
RTK GNSS, IMU, senzorová fúze, Formula Student, LoRa, PPK, Raspberry Pi, CAN, sběr dat
\end{keywordsCZ}

\begin{abstractEN}
This work proposes a system for recording a Formula Student car's
three-dimensional trajectory relative to a local base. Raw GNSS observations,
IMU measurements and a real-time navigation solution are retained for PPK
processing and smoothing. The rover uses a Septentrio mosaic-X5 on a
simpleRTK3B Pro board; the portable base uses a u-blox ZED-F9P on a
simpleRTK2B Budget board. Adafruit ICM-20948 replaces the initially
considered ISM330DHCX for its magnetometer and availability. The design
addresses higher noise, magnetic interference and the chip's end-of-life status.

A 2~GB Raspberry Pi 5 with cooling, M.2 HAT+ and an author-supplied SSD
will run Raspberry Pi OS and read the IMU through SPI and an interrupt.
The proposed accelerometer and gyroscope rate is \SI{562.5}{Hz}; the
magnetometer operates at up to \SI{100}{Hz}. GNSS PPS and GPIO timing
require validation. Two of four XIAO ESP32S3 + Wio-SX1262 kits will form
an 868~MHz RTCM link; two are development/spare units. The base is to
include raw logging and web management. Firmware and transmission require
testing. Two ANN-MB2 antennas need custom mounting. The selected electronics
total about CZK~24,290 excluding VAT, shipping and installation.
A further chapter extends the design to vehicle-wide acquisition of roll/heave,
wheel speed, caliper and tyre temperatures. It adopts a central ADC,
two vehicle CAN interfaces, timing checks and an 8--16 V input. These are outside
the original budget. This is a design proposal, not a report of experiments.
\end{abstractEN}

\begin{keywordsEN}
RTK GNSS, IMU, sensor fusion, Formula Student, LoRa, PPK, Raspberry Pi, CAN, data acquisition
\end{keywordsEN}
